\section{Models Used}

%szb: ez itt egy nagyon hetkikus section lett, több okból:
% - nem szükséges annyi implementation detail - ugyanakkor kell, hogy milyen paraméterekkel finomhangoltál persze (lora alpha, 4-bites kvantálás, stb.)
% - A section cím alapján arra számítottam, hogy itt csak leírod, hogy mely 2 modellt használtad fel, mik ők kb (1-2 mondatban) és ennyi. Ez benne van persze, talán túl részletesen is
% - csak ezeken túl rengetegminden más is itt van, pl. hogy hgoyan finetuneoltál, milyen paramétereket használtál, mivel értékelted ki, stb. Ez így nagyon katyvasz


Two code-specific large language models were selected for the experiments:
\texttt{bigcode/starcoder2-3b} and
\texttt{Qwen/Qwen2.5-Coder-3B-Instruct}. Both models are relatively lightweight
compared to larger code-focused models, making them suitable for training and
evaluation under limited hardware resources, while still providing strong coding
capabilities.

\begin{table}[h]
\centering
\caption{Overview of the evaluated models}
\begin{adjustbox}{max width=\textwidth}
\begin{tabular}{lcccc}
\hline
Model & Parameters & Context Length & Training Tokens & Fine-Tuned \\
\hline
StarCoder2-3B & 3B & 16,384 & 3T+ & Yes \\
Qwen2.5-Coder-3B-Instruct & 3.09B & 32,768 & 5.5T+ & Yes \\
\hline
\end{tabular}
\end{adjustbox}
\label{tab:model_overview}
\end{table}

\subsection{Base Models}

The first set of experiments used the original pre-trained versions of both
models without additional task-specific training.

The \texttt{bigcode/starcoder2-3b} model is a decoder-only Transformer model
with approximately 3 billion parameters. It is part of the StarCoder2 family
and was trained on more than 3 trillion tokens collected from large-scale code
repositories and software-related text sources. The model supports 17
programming languages and provides a context window of up to 16,384 tokens
\cite{starcoder2, starcoder2_hf}. At the time of its release, StarCoder2-3B
outperformed several other code models of similar size on standard programming
benchmarks.

The \texttt{Qwen/Qwen2.5-Coder-3B-Instruct} model is also a decoder-only
Transformer model with approximately 3.09 billion parameters. Unlike
StarCoder2, it is instruction-tuned for coding tasks, which makes it more
suitable for prompt-based code generation and translation. The model supports a
context length of up to 32,768 tokens and was trained on more than 5.5 trillion
tokens. Qwen2.5-Coder models are optimized for programming-related tasks such
as code generation, reasoning, debugging, and repair \cite{qwen_coder,
qwen_coder_report}.

\subsection{Fine-Tuned Models}

In addition to the original pre-trained models, fine-tuned variants of both
StarCoder2 and Qwen2.5-Coder were prepared for the Java-to-C++ translation
task. Fine-tuning was performed using Low-Rank Adaptation (LoRA), which adds
small trainable matrices to selected layers of the model while keeping the
original weights frozen \cite{hu2021lora}. This significantly reduces the
number of trainable parameters and makes fine-tuning more efficient.

Both models were fine-tuned using 4-bit quantization together with LoRA, which
corresponds to the QLoRA approach. In this setup, the original model weights
remain frozen and quantized, while only the LoRA adapter weights are updated
during training. This substantially reduces GPU memory usage and makes it
possible to fine-tune billion-parameter models on consumer-grade hardware. The
models were loaded in 4-bit mode and trained using mixed precision and gradient
checkpointing in order to further reduce memory requirements.

The same LoRA configuration was used for both models. The low-rank dimension
was set to $r = 8$, the LoRA scaling factor to $\alpha = 16$, and the dropout
rate to 0.05. The adapters were applied to both the attention and feed-forward
layers of the models, specifically to the
\texttt{q\_proj}, \texttt{k\_proj}, \texttt{v\_proj},
\texttt{o\_proj}, \texttt{gate\_proj}, \texttt{up\_proj}, and
\texttt{down\_proj} modules. Applying LoRA to both attention and feed-forward
layers is a common approach because it provides stronger adaptation than using
attention layers alone \cite{lora_modules, hu2021lora}.

The maximum sequence length during training was set to 1024 tokens. Both models
were trained for two epochs using a learning rate of $10^{-4}$, cosine learning
rate scheduling, a warmup ratio of 0.03, and gradient accumulation with 16
steps. The effective batch size per device was 1. These settings were selected
to balance training stability and GPU memory consumption.

For both models, the training data was converted into an instruction-style
format in which the model received a Java source code snippet together with a
prompt asking for the corresponding C++ translation. During training, only the
target C++ response tokens contributed to the loss, while the prompt portion
was masked out. This setup allows the models to focus on learning the desired
translation output format rather than memorizing the prompt structure.
